\section{FOSG views, strategies, and kernel games}
\label{sec:fosg-strategies}

The shift from operational machine to game-theoretic structure
happens at this layer.  We define the FOSG view on top of a
machine, exhibit the event-graph FOSG view, build pure and
PMF-behavioural strategy carriers from the view, and present four
canonical kernel games over which the standard
solution-concept inventory is defined.

\subsection{Why FOSG}

The factored-observation stochastic game formalism
\citep{Kovarik2022} indexes strategies by sequences of per-player
observations and packages chance and decision in the same
transition.  Two of its features fit the construction directly.
First, a player's epistemic state is given by their per-step
observation sequence — the same data that the visibility discipline
of \S\ref{sec:source-language} produces from the source program.
Second, FOSG admits a round-batched view of an asynchronous
machine: a single FOSG round collapses many primitive events of
the underlying machine into one transition, which is exactly the
shape we want for the event-graph machine.

\subsection{The FOSG view}
\label{sec:fosg:view}

A \emph{FOSG view} on top of a machine $\Mach$ is a tuple
\[
  \mathcal{V} = \bigl(
    \{\Act_p\}_p,\;
    \activeSet : \State \to \mathrm{Finset}(\Players),\;
    \mathrm{availActions},\;
    \mathrm{transition},\;
    \mathrm{reward}
  \bigr)
\]
of: a per-player round-action alphabet $\Act_p$ (independent of
$\Mach$'s action type); an \emph{active set} $\activeSet(s)$ at each
state; a per-player legal-action predicate
$\mathrm{availActions}(s, p) \subseteq \Act_p$; a transition
kernel
\[
  \mathrm{transition} : (s : \State) \times
    \{\,\sigma : \mathrm{JointAction}\,\Act \mid
      \mathrm{legal}(s, \sigma)\,\}
    \to \PMF{\State},
\]
and a per-player reward function $\mathrm{reward}$ paid on the
transition.  Two side conditions hold: terminal states have empty
active sets, and non-terminal states admit at least one legal joint
action.

\subsubsection*{The event-graph FOSG view}

Concretely, one FOSG round of the event-graph machine
$\Mach_{\EG}$ is the following batched transition.  At a
configuration $C$, the active set is the set of players who own at
least one frontier node.  Each active player simultaneously selects
a write slice for every frontier node they own.  Those player nodes
then execute concurrently, advancing the configuration via the
frontier diamond (\Cref{thm:diamond}).  Every internal node that
becomes enabled in the new configuration then fires
(deterministically for $\mathsf{assign}$ and $\mathsf{reveal}$,
stochastically for $\mathsf{sample}$), and the round ends when the
frontier contains only player-owned nodes again — or, if the
frontier is empty, the round is terminal.  A whole batch of this
kind is one FOSG transition; per-player observations and rewards
are emitted at its boundary.

The FOSG round-action alphabet is therefore static: independently
of the configuration, $\Act_p$ is the type
\[
  \Act_p \;=\; \mathrm{Option}\Bigl(
    (n : \Nodes_p) \to \mathrm{Slice}(n)
  \Bigr),
\]
where $\Nodes_p$ is the set of nodes owned by $p$ in $\EG$ and
$\mathsf{none}$ denotes \emph{not active in this round}.  The
configuration enters at the level of \emph{availability}, not the
alphabet: at configuration $C$, the available actions for $p$ are
those $\mathsf{some}\,a$ whose underlying nodes lie in
$\frontier(C) \cap \{n : \actor(n) = p\}$ and whose slices are
legal there.  Keeping the alphabet static means strategy types do
not depend on the run-time configuration; the per-state menu is
recovered by intersecting with the frontier.
The transition collapses one frontier batch — the player events
owned by the active set, executed concurrently, followed by all
internal events that newly become enabled — into one FOSG step.

\subsubsection*{The stability condition}

The collapse from machine to FOSG is sound only when frontier
choices do not interfere across players.  This is the
\emph{frontier-stability} condition: if action $a$ for player $p$ is
available at $C$, $a$'s node is on $\frontier(C)$ with $a.n \neq
n^\star$, and any other frontier node $n^\star$ executes legally,
then $a$ is still available at the resulting configuration.

\begin{theorem}[Frontier stability for elaborated event graphs]
\label{thm:stability}
For every elaborated graph $\mathsf{eg}(g)$ of a checked program
$g$, frontier stability holds.
\end{theorem}

The proof is a case analysis on the kinds of frontier nodes:
distinct commits do not write to each other's storage (by SSA),
distinct internal nodes likewise, and the
$\mathrm{actionLegal}$ predicate of a commit reads only owner-visible
or public fields, none of which are written by other players'
frontier slices (by the guard-uses-only predicate).

\Cref{thm:stability} is what makes the round-batched view
well-defined: without it, two simultaneously selected frontier
choices could interfere, and the FOSG transition would be
ambiguous.

\subsection{Reachable information states and strategies}
\label{sec:fosg:strategies}

For the bounded FOSG view of a checked program, the set of
\emph{information states} of player $p$ is the set of equivalence
classes of FOSG histories under $p$'s observation.  We work with the
\emph{reachable} subset:
\[
  \mathrm{ReachableInfoState}(p) =
    \bigl\{\, I \in \InfoState_p \mid
      \exists h \in \mathrm{History}.\, h.\mathrm{view}_p = I \,\bigr\},
\]
i.e.\ those information states realised by \emph{some} history in
the FOSG.  Reachability is a purely combinatorial property of
histories, not a property quantified over profiles; this keeps the
strategy carrier definable without circularly referring to the
profiles that range over it.

A \emph{pure strategy} for $p$ is a function
\[
  \sigma^{\mathrm{pure}}_p : \mathrm{ReachableInfoState}(p)
    \to \mathrm{Option}(\Act_p),
\]
returning $\mathsf{none}$ at information states where $p$ is not
active and $\mathsf{some}\,a$ at active ones.  A
\emph{behavioural strategy} is the analogous function valued in
$\PMF{\mathrm{Option}(\Act_p)}$.

The \emph{legal} carriers further restrict: at every reachable
information state where $p$ is active, the chosen $\mathsf{some}\,a$
satisfies $a \in \mathrm{availActions}(h, p)$ for some realising
history $h$.  This is what closes the gap from program-level
\emph{Legal} (existence of a legal action at each commit site) to
\emph{strategic} legality (every strategy in the carrier prescribes
only legal actions): illegal moves are absent from the carrier, not
filtered after the fact.

\designnote{Reachability-pruning matters because illegal moves of
unreachable information states would otherwise pollute the
carrier without affecting any equilibrium.  Pruning is purely
combinatorial: it does not require quantifying over profiles.  The
carrier is still finite by the bounded horizon imposed in
\S\ref{sec:fosg:horizon}.}

\subsection{Bounded horizon and outcome kernels}
\label{sec:fosg:horizon}

A finite event graph admits a finite execution horizon: the source
program has a finite syntax-step count
$h(g) = |\,\mathrm{constructors of }\,g \text{ except }\Ret\,|$, and
no FOSG round can advance the depth counter beyond $h(g)$.  Bounding
the FOSG presentation by $h(g)$ produces a finite-tree-shaped game
to which the GameTheory library's finite-game theorems apply.

The bounded view exposes outcome kernels for both pure and
behavioural profiles:
\begin{align*}
  \mathrm{outcomeFromPure}_h\,\pi      &\in \PMF{\Outcome}, \\
  \mathrm{outcomeFromBehavioral}_h\,\beta &\in \PMF{\Outcome}.
\end{align*}

\subsection{The four canonical kernel games}
\label{sec:fosg:kernel-games}

We can now define the standard strategic-form objects.  A
\emph{kernel game} is the minimal strategic-form carrier on top of
which expected-utility solution concepts make sense:
\begin{equation}
\label{eq:kernel-game}
  \KG \;=\; \bigl(\,
    \{S_p\}_p,\, O,\, K : {\textstyle\prod_p} S_p \to \PMF{O},\,
    u : O \to (\Players \to \mathbb{R})\,\bigr).
\end{equation}
Expected utility is then the obvious functional
\begin{equation}
\label{eq:eu}
  \EU(\Profile, p) \;=\;
  \mathbb{E}_{\omega \sim K(\Profile)}\bigl[u(\omega)(p)\bigr].
\end{equation}

For a checked program $g$, four kernel games are exposed:
\begin{center}
\small
\begin{tabular}{@{}lp{0.55\textwidth}@{}}
\toprule
Name & Carrier and outcomes \\
\midrule
$\mathrm{pureKernelGame}(g)$ &
  pure strategies; outcomes in $\Outcome$ (public payoff). \\
$\mathrm{pmfBehavioralKernelGame}(g)$ &
  PMF behavioural strategies; outcomes in $\Outcome$. \\
$\mathrm{pureBlockedTraceKernelGame}(g)$ &
  pure strategies; outcomes in $\mathrm{BlockedTrace}(g)$. \\
$\mathrm{pmfBehavioralBlockedTraceKernelGame}(g)$ &
  PMF behavioural; outcomes in $\mathrm{BlockedTrace}(g)$. \\
\bottomrule
\end{tabular}
\end{center}

A \emph{blocked trace} is the sequence, indexed by FOSG round, of
all primitive event-block executions that took place during that
round, paired with the configuration reached at the round's end.
It is strictly more informative than the public outcome — which
records only the final public state — and is the natural object on
which to compare the canonical event-graph machine to a backend
implementation, since the latter may produce additional
implementation-specific events within the same round.  The two
axes of the kernel-game table above are then \emph{strategy
carrier} (pure vs.\ PMF-behavioural) and \emph{outcome shape}
(public payoff vs.\ blocked-trace).  The blocked-trace games
\emph{refine} the public games by an outcome projection.

\subsection{Trace-to-public agreement}
\label{sec:fosg:trace-public}

The blocked-trace games are useful for refinement reasoning
(\S\ref{sec:refinement}); the public games are what users
intuitively reason about.  We need to know that they agree, both at
the kernel level and at the level of expected utility.

\begin{theorem}[Trace-to-public agreement]
\label{thm:trace-public}
For every checked program $g$, the public-outcome map applied to
the blocked-trace kernel agrees with the public outcome kernel,
\begin{align*}
  \mathrm{map}(\mathrm{publicOf})\,\circ\,
    K^{\mathrm{pure}}_{\mathrm{blocked}}
  &= K^{\mathrm{pure}}_{\mathrm{public}}, \\
  \mathrm{map}(\mathrm{publicOf})\,\circ\,
    K^{\mathrm{beh}}_{\mathrm{blocked}}
  &= K^{\mathrm{beh}}_{\mathrm{public}},
\end{align*}
and expected utilities are equal under the corresponding profile
inclusions.  Solution concepts (Nash, dominance) transport across
the projection: a strategy is best-response in the public game iff
it is best-response in the blocked-trace game.
\end{theorem}

The proof packages the projection as a kernel-game
\emph{simulation} — a structural morphism that pushes outcome
distributions through a deterministic function while preserving
utility.  In this case the simulation is in fact an
\emph{EU-isomorphism}: the blocked-trace outcome retains
strategically inert per-event-block detail on top of the public
outcome, so the projection is information-losing but not
\emph{utility}-losing, and equilibrium predicates transport in both
directions.  This is the structural reason for the iff in the
trace-to-public statement, and is what distinguishes it from the
backend morphism of \S\ref{sec:refinement}, which is genuinely
one-way: the implementation may admit profiles that no spec
strategy realises.

\subsection{Solution concepts}
\label{sec:fosg:solution}

The standard inventory is defined uniformly on
$\mathrm{pmfBehavioralKernelGame}(g)$.  We list the most important
predicates here and refer to the supplement for the complete table
and the corresponding existence corollaries.

\begin{itemize}\itemsep2pt
\item \textbf{Nash equilibrium.} $\mathrm{IsNash}(\Profile)$ iff
  $\forall p,\, \forall s'_p.\, \EU(\Profile, p) \ge
  \EU(\Profile[p \mapsto s'_p], p)$.
\item \textbf{Best response.} $\mathrm{IsBestResponse}(p, \Profile,
  s)$ iff $s$ maximises $\EU(\Profile[p \mapsto \cdot], p)$.
\item \textbf{Dominance.} $\mathrm{WeaklyDominates}(p, s, t)$ iff
  $s$'s payoff weakly dominates $t$'s for $p$ uniformly in opponents'
  partial profiles.  Strict version with $>$.
\item \textbf{Correlated and coarse correlated equilibrium.}
  $\mathrm{IsCorrelatedEq}(\mu)$~\citep{Aumann1974} for a profile
  distribution $\mu$, and the coarse variant restricting to
  constant deviations.
\item \textbf{Pareto efficiency, social welfare.} $\Profile$ is
  Pareto efficient iff no $\Profile'$ weakly improves all players
  with strict improvement for some $p$.  Social welfare is
  $\sum_p \EU(\Profile, p)$.
\item \textbf{Potential games}~\citep{MonderShapley1996}. A potential
  $\Phi : \prod_p S_p \to \mathbb{R}$ tracking unilateral
  deviations exactly (exact potential) or in sign only (ordinal
  potential).
\item \textbf{Zero-sum and team games.}  Constant total payoff and
  identical per-player utility, respectively.
\end{itemize}

The library exposes finite-game existence corollaries
($\mathrm{mixedNash\_exists}$,
$\mathrm{correlatedEq\_exists}$, $\mathrm{coarseCorrelatedEq\_exists}$)
as wrappers around the GameTheory library's fixed-point and
correlated-equilibrium theorems.

\paragraph{Lean correspondence.}
Supplement~\S6 (machines), \S7.1--\S7.4 (FOSG view, strategies,
kernel games, simulation/bisimulation), \S7.5
(reachable info states), \S7.7 (solution concepts).
